\documentclass[11pt]{article}

\usepackage[margin=1in]{geometry}
\usepackage{amsmath, amssymb}
\usepackage{graphicx}
\usepackage{booktabs}
\usepackage{hyperref}
\usepackage{natbib}

\title{Band Depth Uncertainty Quantification for Chain-of-Thought Reasoning in Large Language Models}
\author{}
\date{}

\begin{document}

\maketitle

% -----------------------------------------------------------------------
\begin{abstract}
Chain-of-Thought (CoT) prompting has become the dominant paradigm for eliciting
reasoning from large language models (LLMs), yet it has created a widening gap
in uncertainty quantification (UQ). When a model generates a step-by-step
reasoning chain, it exposes far richer information than just a final answer---
but existing black-box UQ methods discard this structure entirely, reducing
reliability estimation to counting how often sampled final answers agree.
This is insufficient: two chains that reach the same answer via structurally
divergent reasoning paths are not equally trustworthy, yet answer-level
consistency treats them identically. Approaches that ask the model to verbalize
its own confidence fare no better, yielding scores barely above chance
(AUROC\,$\approx$\,0.50--0.55 in our experiments). The core problem is that
meaningful uncertainty signals are latent in the \emph{geometry} of the
reasoning process---in where chains converge, where they diverge, and how
central each chain is relative to the ensemble---and neither answer-level
counting nor verbal introspection can access that geometry.

We propose \textbf{BD-UQ}, a family of truly black-box UQ methods that operate
directly on the structure of CoT reasoning ensembles using \emph{Band Depth}
(BD), a classical concept from functional data analysis that measures how
centrally a curve lies within a distribution of curves. We lift band depth from
ordered curves to discrete reasoning paths by building a \emph{reasoning graph}
over an ensemble of sampled CoT chains: nodes are semantic clusters of reasoning
steps and edges capture sequential or cross-chain semantic proximity. Three
variants exploit this graph at increasing levels of sophistication---\textbf{pBD}
operates on aligned paths, \textbf{vBD} on individual graph vertices, and
\textbf{gBD} on random walks that explore the full topology of the reasoning
space beyond the specific chains generated. All variants require only a small
sentence encoder and zero additional LLM calls. Extensive experiments across
mathematical reasoning and multi-hop question answering benchmarks show that
gBD consistently and substantially outperforms sampling-based and verbalized
baselines, demonstrating that the reasoning graph is a powerful and previously
untapped source of uncertainty signal.
\end{abstract}

% -----------------------------------------------------------------------
\section{Introduction}
\label{sec:intro}

Large language models (LLMs) have become capable general-purpose reasoners, in
large part due to Chain-of-Thought (CoT) prompting~\citep{wei2022chain}, which
elicits explicit step-by-step reasoning before a final answer and has driven
strong gains on mathematical and multi-hop reasoning tasks~\citep{wang2023selfconsistency, yao2023tree}.
As these models are deployed in increasingly consequential settings, however,
their fluent but sometimes confidently wrong outputs make it essential to know
\emph{when to trust} a given response. Uncertainty quantification (UQ) addresses
this by attaching a confidence score to each generation, ideally one that ranks
correct answers above incorrect ones. The most practically important regime is
\emph{black-box} UQ, where only the generated text is available---no token
logits, hidden states, or gradients---since this is the only access mode offered
by the strongest proprietary models and the cheapest to deploy for open ones.

Existing black-box UQ methods almost universally reduce a response to its
\emph{final answer} and measure agreement across multiple sampled answers.
Self-consistency counts how often the majority answer recurs~\citep{wang2023selfconsistency};
semantic entropy clusters answers by meaning and measures the entropy of the
cluster distribution~\citep{kuhn2023semantic, farquhar2024detecting};
and related approaches compute pairwise answer similarity or
non-contradiction~\citep{lin2023generating, manakul2023selfcheckgpt, chen2024quantifying}.
A second line of work instead asks the model to \emph{verbalize} its own
confidence, either directly or by self-probing in a fresh
context~\citep{xiong2024llms, liu2024uait}. Both families have a common
blind spot: they treat the reasoning chain---the very structure CoT was designed
to produce---as a black box. Answer-level methods discard it entirely, and
verbalized methods rely on the model's unreliable introspection over it.

This blind spot matters because the reasoning chain carries information that the
final answer cannot. Two responses may converge on the same answer through
entirely different reasoning paths---one sound, one a lucky coincidence of
errors---yet answer-level consistency assigns them identical confidence.
Conversely, an ensemble whose chains disagree sharply at intermediate steps
signals genuine model uncertainty even when the final answers happen to match.
Recent work has begun to incorporate reasoning into UQ: CoT-UQ extracts
keywords and importance scores from each step to reweight token
probabilities~\citep{zhang2025cotuq}, and graph-based metrics summarize
relationships among extracted claims~\citep{jiang2024graphuncertainty}. Yet
CoT-UQ still requires token logits (and is therefore not fully black-box) and
operates on a single response rather than an ensemble, while claim-graph methods
target long-form factuality rather than the geometry of reasoning itself. The
question of how to extract uncertainty from the \emph{structure} of an ensemble
of reasoning chains, using text alone, remains open.

We argue that the right signal lives in the \emph{geometry} of the reasoning
ensemble: where chains converge, where they diverge, and how central each chain
is relative to the others. To capture this, we turn to \emph{band depth}, a
concept from functional data analysis that ranks a set of curves from most
central (the functional median) to most outlying~\citep{lopezpintado2009banddepth},
and which has been extended to ensembles of contours and of paths on a
graph~\citep{santer2023contourbanddepth, raj2017pathboxplots}. Building on this
lineage, we embed every reasoning step of every sampled chain with a sentence
encoder~\citep{reimers2019sentencebert} and construct a \emph{reasoning graph}
whose nodes are semantic clusters of steps and whose edges encode sequential and
cross-chain proximity. On this graph we define a family of depth-based
confidence scores---\textbf{pBD}, \textbf{vBD}, and \textbf{gBD}---that measure
how central each chain is within the ensemble. A chain that lies deep inside the
ensemble's reasoning band is trustworthy; an outlying chain is not. Crucially,
all variants are fully black-box and model-free: they require only the sampled
text and a lightweight encoder, with no additional LLM calls.

Our contributions are as follows:
\begin{itemize}
  \item \textbf{Conceptually}, we identify the structure of CoT reasoning
        ensembles as an untapped source of black-box uncertainty, and show that
        the geometric notion of band depth provides a principled way to exploit
        it---going beyond both answer-level consistency and verbalized
        self-probing.
  \item \textbf{Technically}, we introduce a reasoning-graph construction over
        sampled CoT chains and three band-depth confidence estimators on it:
        \textbf{pBD} (path-level, order-sensitive), \textbf{vBD} (vertex-level,
        order-invariant), and \textbf{gBD} (path-level via random walks,
        order-invariant), each requiring no token logits and no extra LLM
        inference.
  \item \textbf{Empirically}, we evaluate on five benchmarks across mathematical
        reasoning~\citep{cobbe2021gsm8k, patel2021svamp, miao2020asdiv} and
        multi-hop question answering~\citep{yang2018hotpotqa, ho2020wikimultihop}
        with two Llama models~\citep{touvron2023llama2, grattafiori2024llama3},
        showing that gBD consistently and substantially outperforms
        sampling-based and verbalized baselines.
\end{itemize}

The remainder of this paper is organized as follows. Section~\ref{sec:prelim}
formalizes the black-box UQ problem and the CoT generation format.
Section~\ref{sec:method} constructs the reasoning graph and defines the pBD,
vBD, and gBD estimators. Section~\ref{sec:related} situates our work within the
UQ and CoT literature. Section~\ref{sec:exp} reports experimental setup,
results, and ablations, and we conclude thereafter.

% -----------------------------------------------------------------------
% \input{sections/preliminary}
% \input{sections/methodology}
% \input{sections/related_works}
% \input{sections/experiment}

% -----------------------------------------------------------------------
\bibliographystyle{plainnat}
\bibliography{references}

\end{document}
